\subsection{HD-D04: explicit sixfold coefficients for the remaining fields}

For \(d=2,11,19,43,67,163\), put \(K=\mathbb Q(\sqrt{-d})\) and
choose the following integral endomorphism:
\[
 a=1+\sqrt{-2}\quad(d=2),\qquad
 a=\frac{1+\sqrt{-d}}2\quad(d>2).
\]
Write \(A=[a]^*\), \(\bar A=[\bar a]^*\) on \(H^6\), and
\(S=A+\bar A\). The actual integral projector used in each
case has the form
\[
 P=p_0+p_1(A+\bar A)+p_2(A^2+\bar A^2)+p_3(A^3+\bar A^3)
   =R e_W.                                               \tag{D04.1}
\]
Every power in this formula is the actual pullback by the
corresponding integral endomorphism \(a^j\) or \(\bar a^j\).
Here are all of its integer coefficients:
\[
\begin{array}{c|r|r|r|r}
d&p_0&p_1&p_2&p_3\\\hline
2&32076&297&-6&-1\\
11&7290&-54&-30&1\\
19&3093750&-5500&-180&1\\
43&6622095018&-1184590&-2310&1\\
67&380746013406&-18787312&-8976&1\\
163&1199676826898478&-4297913560&-132840&1
\end{array}
\]
The positive projector values and the degree bounds for each
signed part are
\[
\begin{array}{c|r|r}
d&R&m\\\hline
2&45056&48114\\
11&26620&46656\\
19&10864656&7000000\\
43&14859858300&11337990400\\
67&607262151936&617921766400\\
163&875011840660800&1854440695686400
\end{array}
\]
For any effective detector of degree \(D\), these operators give
\[
 \boxed{c_t\le C_d\,\delta M_TD^2,}                       \tag{D04.2}
\]
where the exact integer coefficients are
\[
\begin{array}{c|r}
d&C_d\\\hline
2&265503179639\\
11&228308721311\\
19&301570877623313857\\
43&70481797056279525250985339\\
67&1879691667286973076246039776747\\
163&1129506763994798533020373827298661205333
\end{array}
\]
For the order of conductor \(f\), multiply these coefficients
by \(f^{24}\).

\paragraph{Proof of the complete correspondence and constants.}
Set \(N=N_{K/\mathbb Q}(a)\), \(Q=N^3\), and
\[
 s_j=N^j\operatorname{Tr}_{K/\mathbb Q}(a^{6-2j}),
 \quad j=0,1,2,3.
\]
The seven embedding summands of \(H^6\), paired by conjugation,
have these four \(S\)-eigenvalues:
\[
\begin{array}{c|r|r|r|r}
d&s_0&s_1&s_2&s_3\\\hline
2&46&-42&-18&54\\
11&10&21&-45&54\\
19&-54&155&-225&250\\
43&-1638&2189&-2541&2662\\
67&-7326&8687&-9537&9826\\
163&-122958&131159&-136161&137842
\end{array}
\]
All four values are distinct. Choose \(\epsilon=1\) or \(-1\)
so that
\[
 R=\epsilon\prod_{j=1}^3(s_0-s_j)>0.
\]
Then \(\epsilon\prod_{j=1}^3(S-s_j)\) equals \(R\) on the
two Weil lines and zero on every other embedding summand.
This proves the projector identity on all of \(H^6\).

The exact cohomological relations
\[
 A\bar A=[N]^*=Q^2,\quad
 S^2=A^2+\bar A^2+2Q^2,\quad
 S^3=A^3+\bar A^3+3Q^2S
\]
expand that polynomial to (D04.1), with the displayed
coefficients. These relations are used to verify the
cohomology class of the explicit integral cycle operator;
no equality in Chow groups is presumed.

Since \(s_3=2Q\), the signed degree sum is
\(p_0+2p_1Q+2p_2Q^2+2p_3Q^3=0\).
Thus the positive and negative effective parts each have
degree at most \(mD\), where \(m\) is the positive weighted
coefficient sum. The table evaluates it exactly.
The integers \(p_3=\pm1\) show that these displayed operator
coefficients admit no further common integral divisor.

In every row \(r=s_0=\operatorname{Tr}(a^6)\) is even.
Use the smaller integral rotation
\[
 Z=A-r/2,\qquad
 s=Q^2-r^2/4,\qquad L=Q+|r|/2.
\]
On \(W\) it satisfies \(Z^\dagger=-Z\) and \(Z^2=-s\);
its signed degree cost is at most \(L\).
The respective pairs \((s,L)\) are
\[
 (200,50),\ (704,32),\ (14896,152),\
 (1100800,2150),\ (10720000,8576),\
 (970436800,130400).
\]
The proof of HD-D03 applies with these exact data. Its
multiplier bound is
\(\lceil sR^2/36\rceil\), and its signed-degree coefficient is
\[
 C_d=\left\lceil\frac{mR\sqrt{s(s+L^2)}}6\right\rceil.
\]
For every row the latter is at least the former.
The accompanying exact script evaluates this ceiling as the
least nonnegative integer \(k\) with
\[
 (6k)^2\ge(mR)^2s(s+L^2),
\]
and stores the strict inequality for \(k-1\) as well.
All seven projector eigenvalues and both effective signed
degree sums are checked in the integer quadratic ring.

Finally, \(a/\bar a=a^2/N\) has a nonintegral coefficient in
the displayed maximal-order integral basis: for \(d=2\) the
coordinates are \((-1/3,2/3)\), and for \(d>2\) they are
\((-1,1/N)\). It cannot be a root of unity.
For the order of conductor \(f\), replace every term
\([u]^*\) of (D04.1) and of \(Z\) by \([fu]^*\).
All are actual integral endomorphisms, while on \(H^6\)
they multiply the corresponding operator by \(f^6\).
Thus \(R,m\) gain \(f^6\), \(s\) gains \(f^{12}\), and
\(L\) gains \(f^6\). The quadratic extraction coefficient
gains \(f^{24}\), proving the conductor assertion.
\(\square\)

The exact reproducible files are
\texttt{remaining\_field\_projectors.py} and its JSON.
They contain every expanded endomorphism term, all field
coordinates, all signed masses and both integer ceiling witnesses.

\paragraph{The fourfold extension of HD-D04.}
Use the same field elements on \(H^4\).  In this paragraph put
\(Q=N(a)^2\), so that \(A\bar A=Q^2\) on \(H^4\).
The three conjugate-paired \(S\)-eigenvalues are
\[
 s_j=N(a)^j\operatorname{Tr}(a^{4-2j}),\qquad j=0,1,2.
\]
The complete numerical data for these eigenvalues and the actual
five-term cycle operator
\[
 P=p_0+p_1(A+\bar A)+p_2(A^2+\bar A^2)=R e_W
                                                               \tag{D04.3}
\]
are as follows:
\[
\begin{array}{c|r|r|r|r|r|r}
d&s_0&s_1&s_2&p_0&p_1&p_2\\\hline
2&-14&-6&18&54&-12&1\\
11&7&-15&18&108&3&-1\\
19&31&-45&50&1000&5&-1\\
43&199&-231&242&26620&11&-1\\
67&511&-561&578&157216&17&-1\\
163&3199&-3321&3362&5513680&41&-1
\end{array}
\]
For \(d=2\), the projector is \((S+6)(S-18)\), reduced using
\(S^2=A^2+\bar A^2+2Q^2\); its value on the Weil plane is
\(256\).  For every odd \(d\) in the table, write
\(b=(d+1)/4=N(a)\).  Then
\[
 (s_0,s_1,s_2)=(1-4b+2b^2,\ b-2b^2,\ 2b^2),
\]
and the projector is
\[
 -(S-s_1)(S-s_2)
 =2b^3(b-1)+b(A+\bar A)-(A^2+\bar A^2).
\]
The identities
\(s_0-s_1=(b-1)(4b-1)=(b-1)d\) and
\(s_0-s_2=-d\) give \(R=(b-1)d^2>0\).
Thus (D04.3) has the correct value on both Weil lines and vanishes
on all three other embedding summands.  Its signed effective parts
have equal degree bound \(mD\), where \(m=216\) for \(d=2\),
and \(m=2b^4\) for odd \(d\).  These assertions use the displayed
actual endomorphism sum; the polynomial reductions assert an
identity on cohomology only.

For \(d=2\), take \(Z=A+7\), yielding \((s,L)=(32,16)\).
For odd \(d\), \(s_0\) is odd, so take the integral rotation
\[
 Z=2A-s_0,\qquad
 s=4Q^2-s_0^2=(2b-1)^2d,\qquad
 L=2Q+|s_0|=(2b-1)^2.
\]
The equality for \(L\) holds because \(b\ge3\) in these rows.
The exact extraction constants are
\[
\begin{array}{c|r|r|r|r}
d&R&m&s&L\\\hline
2&256&216&32&16\\
11&242&162&275&25\\
19&1444&1250&1539&81\\
43&18490&29282&18963&441\\
67&71824&167042&72963&1089\\
163&1062760&5651522&1069443&6561
\end{array}
\]
\[
\begin{array}{c|r|r}
d&A_0&A_1\\\hline
2&524288&2654208\\
11&4026275&9751872\\
19&802256076&3186464101\\
43&1620768084075&17222775868991\\
67&94098319457472&1818065274326025\\
163&301972956567649200&20627581863862037373
\end{array}
\]
Here
\[
 A_0=\left\lceil\frac{sR^2}{4}\right\rceil,\qquad
 A_1=\left\lceil\frac{mR\sqrt{s(s+L^2)}}2\right\rceil.
\]
For odd \(d\), the second expression is also the exact ceiling of
\[
 2b^5(b-1)(2b-1)^2d^2\sqrt d.
\]
Every table entry satisfies \(A_1\ge A_0\), and therefore a
fourfold detector of degree \(D\) gives
\[
 \boxed{c_t\le A_1\,\delta M_TD^2.}                       \tag{D04.4}
\]
The certificate evaluates \(A_1\) as the least integer \(k\)
such that \((2k)^2\ge(mR)^2s(s+L^2)\), and records the strict
inequality for \(k-1\); no floating-point ceiling is used.
It also tests all five embedding summands and each signed mass.
Replacing every displayed term \([u]^*\) by \([fu]^*\) multiplies
\(R,m,L\) by \(f^4\) and \(s\) by \(f^8\), giving the exact
conductor extraction factor \(f^{16}\).
\(\square\)

